\documentclass[11pt,a4paper]{article}

% --- PAQUETS DE BASE ---
\usepackage[utf8]{inputenc}
\usepackage[T1]{fontenc}
\usepackage[french]{babel}
\usepackage{amsmath, amsfonts, amssymb}
\usepackage{geometry}
\geometry{top=2cm, bottom=2.5cm, left=2cm, right=2cm}

% --- DESIGN ---
\usepackage{xcolor}
\definecolor{BrandPrimary}{HTML}{0D9488}
\definecolor{BrandDark}{HTML}{111827}

\usepackage{tcolorbox}
\tcbuselibrary{skins, breakable}

\usepackage{titlesec}
\titleformat{\section}{\color{BrandPrimary}\normalfont\Large\bfseries}{\thesection}{1em}{}[{\titlerule[0.8pt]}]

% --- POLICE ---
\usepackage{helvet}
\renewcommand{\familydefault}{\sfdefault}

\begin{document}

\begin{tcolorbox}[enhanced, colback=BrandDark, colframe=BrandPrimary, arc=10pt, boxrule=2pt, halign=center]
    \vspace{0.3cm}
    {\Huge \bfseries \color{white} ZENITH LEARN} \\
    \vspace{0.2cm}
    {\Large \color{BrandPrimary} IA Éthique : Modélisation TWIST 04 (V2 - Strict Math)} \\
    \vspace{0.3cm}
\end{tcolorbox}

\section{Modélisation de l'Indice de Potentiel ($SPI$)}
Le TWIST 04 introduit la décorrélation entre complétion et action. Mathématiquement, nous passons d'une fonction linéaire simple à un modèle de pondération par "rupture de domaine".

\subsection{La Formule du Success Potential Index}
Le score de potentiel entrepreneurial ($SPI$) est la somme pondérée des vecteurs d'action et des vecteurs académiques.
\begin{equation}
SPI(u) = \underbrace{w_A \cdot \frac{A_u}{A_{max}} + w_B \cdot B_u}_{\text{Domaine Action}} + \underbrace{w_E \cdot S_E + w_C \cdot S_C^{adj}}_{\text{Domaine Académique}}
\end{equation}

\subsection{Explication détaillée des termes}
\begin{itemize}
    \item \textbf{$w_A$ (Poids de l'Action Terrain) :} Fixé à 0.3. Il sanctionne ou récompense la mise en pratique.
    \item \textbf{$A_u$ (Somme des actions terrain accomplies) :} Variable discrète (nombre d'actions valides).
    \item \textbf{$A_{max}$ (Seuil de référence) :} Fixé à 5. Indique qu'au-delà de 5 actions terrain, le score sature pour ce terme (évite le biais de quantité).
    \item \textbf{$w_B$ (Poids du Lancement d'Entreprise) :} Fixé à 0.4. C'est le \textbf{terme dominant} de l'équation. Le lancement d'une entreprise est le signal de succès ultime.
    \item \textbf{$B_u$ (Variable binaire) :} 1 si entreprise lancée, 0 sinon.
    \item \textbf{$w_E$ (Engagement) et $w_C$ (Complétion) :} Fixés à 0.15 chacun. Ils représentent le "Audit de Soutien". 
    \begin{itemize}
        \item \textbf{Concrètement :} Même si un utilisateur a 100\% en académique ($S_C^{adj}=1, S_E=1$), son score global sans action sera plafonné à $0.15 + 0.15 = 0.30$. 
        \item À l'inverse, un utilisateur avec 20\% académique mais une entreprise lancée aura $SPI > 0.40$.
    \end{itemize}
\end{itemize}

\section{Impact Dataset}
\begin{tabular}{|l|p{10cm}|}
\hline
\rowcolor{BrandDark} \color{white} Colonne & \rowcolor{BrandDark} \color{white} Transformation \\ \hline
\texttt{action\_completed} & Devient le numérateur de l'indice terrain. \\ \hline
\texttt{business\_launched} & Devient l'activateur du terme dominant ($w_B$). \\ \hline
\end{tabular}

\end{document}
